% =====================================================================
%  夏ゼミ配布資料（レジュメ） 松本 鋭
%  コンパイル: uplatex natsuzemi_resume_matsumoto.tex && dvipdfmx natsuzemi_resume_matsumoto.dvi
%  （platex でも可。lualatex を使う場合は documentclass を ltjsarticle に変える）
%  ※ 日付・所属・掲載する数値は、レジュメ作成講座の指示に合わせて直す前提
% =====================================================================
\documentclass[twocolumn,a4paper,9pt]{jsarticle}

\usepackage{amsmath,amssymb}
\usepackage{bm}
\usepackage{graphicx}
\usepackage{url}
\usepackage[top=20mm,bottom=20mm,left=17mm,right=17mm]{geometry}

\setlength{\columnsep}{7mm}
\renewcommand{\baselinestretch}{0.98}

% 表題は手組みにする。jsarticle の \maketitle は二段組だと内部で \twocolumn を
% 呼ぶため、\twocolumn[...] の中に入れると入れ子になって壊れる
% （Argument of \@newline has an extra } の原因）。

\begin{document}

\twocolumn[%
  \noindent\hfill{\small 2026/09/XX 夏ゼミ配布資料}\par
  \vskip 3mm
  \begin{center}
    {\LARGE\bfseries SELDモデルを用いた\\
     難聴歩行者向け屋外危険音通知システムの構築\par}
    \vskip 3mm
    {\large 松本 鋭\footnotemark[1]\par}
    \vskip 1mm
    {\normalsize Satoshi Matsumoto\par}
  \end{center}
  \vskip 5mm
]
\footnotetext[1]{東京理科大学 創域理工学部 情報計算科学科 4年}

% ---------------------------------------------------------------
\section{研究背景と目的}

歩行者は，接近する車やサイレン，踏切といった危険の察知を聴覚に大きく依存している．
難聴者はこの情報が欠落するため，屋外歩行における事故リスクが高い．
この課題はイヤホンを装着した歩行者や高齢者にも共通する．

一方で「機械の耳」は，疲労せず全方位を常時監視でき，
センサが対応すれば人間の可聴域外の信号も扱える．
したがって，単なる聴覚の代替にとどまらない支援が期待できる．

本研究の目的は，屋外の危険音8クラスについて
\textbf{種類・方向・距離}を同時に推定し，
それを\textbf{危険度3段階の通知}に変換して届けるシステムを構築することである．
検出対象はサイレン・クラクション・バック音・自転車ベル・車（EV・大型含む）・
踏切と列車・キックボード・バイクの8クラスとし，
将来的には首元の振動で伝える構成を想定する．行動の判断は利用者に委ねる．

% ---------------------------------------------------------------
\section{基礎知識}

\subsection{SELD}
SELD（Sound Event Localization and Detection）は，
画像における物体検出の音・方向版にあたるタスクであり，
音響イベント検出（SED，何がいつ鳴っているか）と
音源定位（SSL，どの方向から鳴っているか）を，
複数の音源について同時に解く\cite{adavanne}．

\subsection{音源距離推定の拡張}
本研究では SELD に音源距離推定（SDE）を加え，
\textbf{種類・方向・距離}の同時推定へ拡張する．
距離は通知の危険度を決める量であり，
「自分に近づいてくる車」と「遠方を通り過ぎる車」を区別するために不可欠である．

% ---------------------------------------------------------------
\section{関連研究}

\subsection{PSELDNets}
Hu ら\cite{pseldnets}は，1,167時間・170クラスの大規模合成データで事前学習した
SELD の基盤モデルを提案し，主要データセットで最高性能を達成した．
FOA形式の4ch空間音響を入力とし，
同一クラスの音源を最大3つまで同時に検出できる
multi-ACCDOA\cite{maccdoa} 出力を持つ．

\subsection{DynamicSound}
Barbisan ら\cite{dynamicsound}は，室内残響ツールでは扱えなかった
屋外・移動音源の状況を，伝搬遅延・ドップラー・距離減衰・大気吸収・地面反射
という物理で明示的にモデル化した屋外音響シミュレータを公開している．

\subsection{ウェアラブル音響支援}
歩行者安全を対象としたウェアラブル音響システム\cite{paws}や，
頭部装着デバイスによる SELD データセット\cite{wearable}は存在する．
しかし，\textbf{歩行者視点で複数の危険クラスを距離まで含めて扱った研究}は，
調査した範囲では見当たらない．この空白が本研究の位置づけである．

% ---------------------------------------------------------------
\section{提案手法}

\subsection{合成データの自作}
屋外 SELD の最大の障害は，本研究の条件
（屋外・歩行者視点・距離ラベル付き）を満たす学習データが存在しないことである．
そこで物理シミュレーションを自作し，4ch の音と
方向・距離のラベルを同一の計算経路から同時に出力する．
音源の音量・周波数は日本の法規・省令・実測研究に出典を持たせて層別に管理した．
合成にはもう一つ利点があり，物理を自分で保持しているため，
後から要素を一つずつ外す ablation が可能な土俵になる．

学習は PSELDNets を距離ヘッドごとファインチューニングし，
自作の合成10,200クリップ（8クラス・約28時間）を用いた．

\subsection{2層構造}
システムは知覚層と通知層の2層からなる．
知覚層は8クラスの検出・方向・距離を0.1秒ごとに推定する．
通知層はその推定から「いつ・どの強さで」伝えるかを決める．
危険のたびに全て鳴らすと実用に耐えないため，
\textbf{通知の頻度そのものを設計対象}としている点が本システムの特徴である．

\subsection{通知層の改良：最接近距離の予測}
当初の通知層は推定距離のしきい値のみで判定していた．
これは速度を見ていないため，相手が速いほど余裕が縮むという構造的な弱点を持つ．
実測でも，至近警告が出てから最接近までの余裕は中央値0.00秒であった．

そこで，距離の縮み方 $\dot d$ と方位の振れ方 $\dot a$ から
\textbf{最接近の距離と時刻}を等速直線運動として解き，これで判定する方式に改めた．
相対位置 $\bm r$（大きさ $d$）と相対速度 $\bm v$ を
動径成分と接線成分に分けると
%
\begin{equation}
  |\bm v|^{2}=\dot d^{2}+(d\dot a)^{2},\quad
  t_{\mathrm{cpa}}=\frac{-d\,\dot d}{|\bm v|^{2}},\quad
  d_{\mathrm{cpa}}=\frac{d^{2}|\dot a|}{|\bm v|}
\end{equation}
%
となる．予測最接近距離 $d_{\mathrm{cpa}}$ が小さく，
かつ到達時刻 $t_{\mathrm{cpa}}$ が短いときに警告する．
これは「方位が変わらないまま近づく物体は衝突コースにある」という
船舶の見張りの原則を，距離推定を得たことで定量化したものにあたる．
これにより，正面から接近する対象には遠方で鳴り，
横を通り過ぎるだけの対象には鳴らない，という出し分けが可能になる．

なお，接近速度のみを用いる到達時間（TTC）による判定も試したが，
横を通り過ぎるだけの対象にも警告が出るため，
安全な対象を鳴らさずに済んだ割合が82.1\%から15.7\%へ崩壊した．
接近しているかどうかだけでは不十分であり，
進行方向の情報が不可欠であることが確認できた．

% ---------------------------------------------------------------
\section{評価と結果}

\subsection{評価手順}
基準を先に決め，学習にも検証にも用いていない
同一設計・新乱数の1,800クリップを新たに生成し，
\textbf{1回だけ}採点する手順を事前登録した．
検証用データはモデル選択と通知しきい値の決定にのみ用いている．

\subsection{知覚層}
8クラスの検出率は 97.2--100\%，位置込みの総合 F 値は 91.3\% であった．
方向誤差の中央値は 1--5 度である
（出力・正解ともに1度刻みで記録しているため，これ以上は分解できない）．
距離は，通知が作動する至近5m以内で中央絶対誤差 0.21m であり，
相対誤差でみると 1.5--30m の範囲でおおむね 6--9\% と安定している．
エンジン音を持たない静音 EV のみの場面でも検出率は 98.99\% であった．

\subsection{通知層}
1.5m 以内まで接近する危険な車のうち，至近警告が届いた割合は
距離しきい値のみの方式で 69.9\% であった．
同一の採点器・同一の分母のまま最接近予測の方式に差し替えると
\textbf{85.0\%} へ向上し，
警告から最接近までの余裕の中央値は 0.00 秒から \textbf{0.80 秒}に伸びた．
一方で，安全な車を鳴らさずに済んだ割合は 90.4\% から 85.2\% へ低下する．
早く鳴らすことと余計に鳴らさないことは両立せず，
この対価は正直に報告すべき性質のものである．

% ---------------------------------------------------------------
\section{今後の方針}

\begin{itemize}\itemsep=0pt
  \item \textbf{実環境録音による検証}：4ch FOA マイクを胸部に装着し，
        統制録音・機会捕捉・負例の3層で計100本を収録する．
        加えて，装着者自身の歩行が推定に与える影響を測るため，
        同一条件を静止と歩行で対にした100本（50組）を別枠で収録する（計200本）．
        車の至近帯は安全上再現できないため，
        自転車とキックボードによる部分検証にとどめることを先に宣言する．
  \item \textbf{物理 ablation}：距離減衰・ドップラー・地面反射・大気吸収を
        1つずつ外した4条件を同一の土俵で採点し，
        「合成データのどの物理が実環境性能に必要か」に直接答える．
  \item \textbf{距離推定の精度向上}：入力を正解値に置き換えた上限実験から，
        通知の余裕時間を伸ばす鍵が距離推定の精度と複数音源の追跡にあることが分かっている．
        これを次の目標に加える．
\end{itemize}

% ---------------------------------------------------------------
\begin{thebibliography}{9}\small
\bibitem{pseldnets} J.~Hu \textit{et al.},
  ``PSELDNets: Pre-trained Neural Networks on a Large-scale Synthetic Dataset
  for Sound Event Localization and Detection,'' IEEE/ACM TASLP, 2024.
\bibitem{dynamicsound} L.~Barbisan, M.~Levorato, F.~Riente,
  ``DynamicSound simulator for simulating moving sources and microphone arrays,''
  arXiv:2601.15433, 2026.
\bibitem{maccdoa} K.~Shimada \textit{et al.},
  ``Multi-ACCDOA: Localizing and Detecting Overlapping Sounds from the Same Class
  with Auxiliary Duplicating Permutation Invariant Training,'' ICASSP, 2022.
\bibitem{adavanne} S.~Adavanne, A.~Politis, J.~Nikunen, T.~Virtanen,
  ``Sound Event Localization and Detection of Overlapping Sources Using
  Convolutional Recurrent Neural Networks,'' IEEE J.~Sel.~Top.~Signal Process.,
  vol.~13, no.~1, pp.~34--48, 2019.
\bibitem{mesaros} A.~Mesaros \textit{et al.},
  ``Joint Measurement of Localization and Detection of Sound Events,''
  WASPAA, pp.~333--337, 2019.
\bibitem{paws} D.~de~Godoy \textit{et al.},
  ``PAWS: A Wearable Acoustic System for Pedestrian Safety,''
  ACM/IEEE IoTDI, 2018.
\bibitem{wearable} K.~Nagatomo, M.~Yasuda, K.~Yatabe, S.~Saito, Y.~Oikawa,
  ``Wearable SELD dataset: Dataset for sound event localization and detection
  using wearable devices around head,'' ICASSP, 2022.
\bibitem{findlater} L.~Findlater \textit{et al.},
  ``Deaf and Hard-of-hearing Individuals' Preferences for Wearable and Mobile
  Sound Awareness Technologies,'' CHI, 2019.
\bibitem{vogel} K.~Vogel,
  ``A comparison of headway and time to collision as safety indicators,''
  Accident Analysis \& Prevention, vol.~35, no.~3, pp.~427--433, 2003.
\end{thebibliography}

\end{document}
